% India and Hyderabad slides (all numbers from numbers_india.tex)
\begin{frame}{India: what the nine-crop index says}
\begin{columns}[T]
\column{0.52\textwidth}\small
\textbf{Ranking for India:} \indOrder
\begin{itemize}
\item Wheat \indWheatFFour: \indWheatShIron\% of national iron, \indWheatShKcal\% of energy
\item Wheat today: mostly water; its \emph{future} rise is almost all heat
\item Chickpea \indChickpeaFFour: all water --- the highest rainfed deficit in the panel
\item 10\% wheat loss: energy \indWheatShockKcal\%, iron \indWheatShockIron\%, MAR \indWheatShockMar\ pts
\end{itemize}
\column{0.48\textwidth}\small
\textbf{A concentrated food supply}
\begin{itemize}
\item Effective food items supplying energy: India \divIndKcal, China \divChnKcal, USA \divUsaKcal, world \divWorldKcal
\item Cereals: \divIndCereal\% of energy
\item Wheat: largest iron source (\divIndIronTopShare\%)
\end{itemize}
\end{columns}
\src{FAO food balance sheets 2021--23 (Shannon diversity of supply, after Remans et al.\ 2014).}
\end{frame}

\begin{frame}{Hyderabad: a rice diet, unlike the national wheat story}
\begin{columns}[T]
\column{0.5\textwidth}\small
\textbf{Urban Telangana} (nearest official data to Hyderabad), kg per person per month
\begin{itemize}
\item Rice \hqTgUrbanKgRice\ vs wheat \hqTgUrbanKgWheat\ ($\times$\hqTgUrbanRiceToWheat; urban India $\times$\hqInUrbanRiceToWheat)
\item Cereals: \hqKcalCerealTgUrban\% of energy, \hqProtCerealTgUrban\% of protein
\item Protein-source diversity \hqProtEffTgUrban\ of 5 groups (urban India \hqProtEffInUrban)
\item Animal protein \hqProtAnimalTgUrban\% (India urban \hqProtAnimalInUrban\%)
\end{itemize}
\column{0.5\textwidth}\small
\textbf{Food groups, Telangana women: at least weekly}
\begin{itemize}
\item Pulses \hqWeeklyPulses\%, milk \hqWeeklyMilk\%, eggs \hqWeeklyEgg\%, meat \hqWeeklyMeat\%
\item Leafy vegetables \hqWeeklyLeafy\%, fruit \hqWeeklyFruit\%, fish \hqWeeklyFish\%
\end{itemize}
\textbf{Gaps:} fruit, leafy vegetables, fish
\end{columns}
\src{MoSPI HCES 2023--24 and NSS 594 (state, rural/urban); NFHS-5 Telangana state report (2019--21).}
\end{frame}

\begin{frame}{Nutrient diversity: where Hyderabad's iron comes from}
\begin{columns}[T]
\column{0.55\textwidth}\small
Share of iron from listed grains and pulses (HCES quantity $\times$ IFCT 2017):
\begin{itemize}
\item Rice \hqTgUrbanFeRice\% (urban India \hqInUrbanFeRice\%)
\item Wheat \hqTgUrbanFeWheat\% (urban India \hqInUrbanFeWheat\%)
\item Chickpea \hqTgUrbanFeChickpea\%, lentil \hqTgUrbanFeLentil\%, tur \hqTgUrbanFeTur\%, moong \hqTgUrbanFeMoong\%, groundnut \hqTgUrbanFeGroundnut\%
\end{itemize}
Rice: \SI{0.65}{mg} iron per \SI{100}{g}; chickpea and lentil dal 6.1--7.1.
\column{0.45\textwidth}\small
\textbf{Does the health data agree?} (NFHS-5, Hyderabad district)
\begin{itemize}
\item Anaemic: children \hydAnaemChild\%, non-pregnant women \hydAnaemNPW\%
\item Underweight women only \hydBmiLow\%; overweight/obese \hydOverweight\%
\end{itemize}
$\Rightarrow$ energy is adequate; \textbf{iron is the problem} --- the index's finding, locally.
\end{columns}
\src{Telangana, CNNS 2016--18: low ferritin \cnnsIronCFour\% (1--4 y), \cnnsIronAdol\% (adolescents); zinc deficiency \cnnsZincAdol\% (adolescents).}
\end{frame}

\begin{frame}{How would climate affect Hyderabad's diet?}
\begin{columns}[T]
\column{0.5\textwidth}\small
\textbf{Local climate} (grid cell; MRI-ESM2-0, SSP2-4.5)
\begin{itemize}
\item Mean $T_{\max}$ \hydTmaxBase\,$\to$\,\hydTmaxRecent\,\si{\celsius} (1971--2000 $\to$ 2001--24); rain \hydRainBase\,$\to$\,\hydRainRecent\,mm
\item 2041--60: April--May $T_{\max}$ $+$\hydAMFutChange\,\si{\celsius} (\hydAMFutLevel); rain \hydRainFutRatio\%
\end{itemize}
\textbf{Telangana crops}
\begin{itemize}
\item Rice: \tgAreaRice\% of cropland, \tgRiceRainfed\% rainfed $\Rightarrow$ almost no modelled stress
\item Paddy \tgPaddyAreaMha\,Mha (2023--24) vs \tgPaddySpamMha\,Mha in the 2010 map: $\times$\tgPaddyGrowth
\end{itemize}
\column{0.5\textwidth}\small
\textbf{Effect on iron from staples}
\begin{itemize}
\item Climate-driven yield change: rice \hqdYRice\%, chickpea \hqdYChickpea\%, wheat \hqdYWheat\% (national)
\item Iron from these foods: \hqIronChangeUrban\%
\item 10\% loss of rice: iron \hqIronTenRiceUrban\%; of wheat \hqIronTenWheatUrban\%
\end{itemize}
\begin{block}{Reading}Small next to the existing gap: Hyderabad's risk is structural; the climate route is irrigation, rainfed pulses and imported wheat.\end{block}
\end{columns}
\src{Conditional screening estimates; no CO$_2$ fertilisation, adaptation or market response.}
\end{frame}
